\mxunit{terminale-demonstrations}{Répertoire des démonstrations}
Ce répertoire indique où retrouver les preuves de cours. Il ne fixe pas un nombre universel de « démonstrations exigibles » : les regroupements varient selon la présentation.
\begin{longtable}{@{}p{.69\textwidth}p{.25\textwidth}@{}}\toprule\textbf{Résultat et démarche}&\textbf{Dossier}\\\midrule\endhead
Récurrence : forme explicite, somme, divisibilité&Ch. 1, p. \pageref{ch:1}\\
Suite croissante non majorée ; comparaison vers l'infini&Ch. 2, p. \pageref{ch:2}\\
Bernoulli et limites de $q^n$&Ch. 2, p. \pageref{ch:2}\\
Limites de l'exponentielle et croissance comparée&Ch. 3, p. \pageref{ch:3}\\
Courbe convexe au-dessus de ses tangentes&Ch. 4, p. \pageref{ch:4}\\
Dérivée de ln ; limite de $x\ln x$ en $0^+$&Ch. 6, p. \pageref{ch:6}\\
Deux primitives diffèrent d'une constante ; $y'=ay$&Ch. 7, p. \pageref{ch:7}\\
Primitive d'une aire (cas croissant) ; intégration par parties&Ch. 8, p. \pageref{ch:8}\\
Pascal par deux méthodes ; nombre de parties&Ch. 9, p. \pageref{ch:9}\\
Projeté orthogonal et distance minimale&Ch. 11, p. \pageref{ch:11}\\
Équation cartésienne issue d'un vecteur normal&Ch. 12, p. \pageref{ch:12}\\
Expression de la loi binomiale&Ch. 13, p. \pageref{ch:13}\\
Espérance et variance binomiales ; concentration et LGN&Ch. 14, p. \pageref{ch:14}\\\bottomrule\end{longtable}
